% Nouveau EGA V (cote 25): archive cote 25, site page 141 (PDF page 142).
% English translation of the checked French annotated master.
% AI-translated by GPT-6 Astra (OpenAI, ChatGPT Pro mode).
\scanpage{25/141}
\mspage{\hand{2}}

\uline{Remarks }4. \var{\inserted[typed]{a)}}{a)}{insertion}{ega25_25_141-v0001} From the \var{\xout{xx}}{}{typewriter strike-out}{ega25_25_141-v0002} criterion that will be developed immediately below, it follows
that in Definition 3 it suffices to test with a single point
with values in a k'---one can, for example, \unsure{take} $k'=k(x)$ or $\hand{\overline{\textcolor{black}{k(x)}}}$, and
the canonical point with values in this k'. b) It follows from Remark
2 that x is geometrically nonsingular for $\hand{\phi}$ iff $\hand{\phi}_x \hand{\neq} 0$ and
$V(\hand{\phi})$ is \uline{smooth} over k at x. c) Suppose that we have \var{\inserted[typed]{a}}{a}{insertion}{ega25_25_141-v0004} prescheme X smooth
over another \unsure{S|Y}, a section $\hand{\phi}$ of $\underline{O}_X$ and an $x \hand{\in} V(\hand{\phi})$; then we say that x
is a singular (resp. supersingular) zero \uline{relative to \unsure{S|Y}} if it is so
relative to k(s) on the fiber $X_s$ (s being the image of x in \unsure{S|Y}). d) Under
the conditions of Def. 1, one sees immediately that the singularity, resp.
\var{\margin{$\|$}}{\placedmargin{$\|$}}{marginal note placed}{ega25_25_141-v0006} supersingularity, of an $x \hand{\in} V(\hand{\phi})$ for $\hand{\phi}$ is unchanged when one
replaces $\hand{\phi}$ by $\hand{\phi}'=u\hand{\phi}$, where u is a unit at x. It follows at once that
\var{\xout{the notion}}{}{typewriter strike-out}{ega25_25_141-v0007} Definition 1, and consequently also Def. 3, extend in an
obvious way to the case where $\hand{\phi}$ is a section of an invertible Module $\underline{L}$ (so as
to recover the initial definitions for $L = \underline{O}_X$).

\var{\xout{Consider the section xxxxx d$^2$}}{}{typewriter strike-out}{ega25_25_141-v0008} Let X be a prescheme smooth over
another Y, and let $\hand{\phi}$ be a section of $\underline{O}_X$, giving a section $d^2_{X/Y}\hand{\phi}$ of
$P^2_{X/Y}$, reducing to the section $d^1_{X/Y}\hand{\phi}$ of $P^1_{X/Y}$, which itself
% From Herewith to the page end, the source is framed, with small ticks along its lines.
reduces to the section $d^0\hand{\phi}$ \var{\inserted[typed]{$=\hand{\phi}$}}{$=\hand{\phi}$}{insertion}{ega25_25_141-v0009} of $P^0_{X/Y} = \underline{O}_X$. With this understood, we have
\uline{Proposition} 5 The set of zeros of $d^{\hand{0}}\hand{\phi}$ (resp. $d^1\hand{\phi}$ \hand{)} \var{\deleted{resp. $d^2\hand{\phi}$)}}{}{deletion}{ega25_25_141-v0010}
\var{\margin{retain the text opposite}}{\placedmargin{retain the text opposite}}{marginal note placed}{ega25_25_141-v0011} is respectively equal to the set $V(\hand{\phi})$ of zeros of $\hand{\phi}$ (resp. to
the set $V(\hand{\phi})^{\mathrm{sing}}$ of zeros of $\hand{\phi}$ singular relative to S\hand{)} \var{\deleted{resp. to the set}}{}{deletion}{ega25_25_141-v0012}
\var{\deleted{$V(\hand{\phi})^{\mathrm{sup\,sing}}$ of zeros of $\hand{\phi}$ supersingular relative to S)}}{}{deletion}{ega25_25_141-v0013}.

The first assertion is trivial; the second is simply
the Jacobian criterion, or, if one prefers, follows from the canonical
isomorphism $\underline{m}_x/\underline{m}_x^2 \hand{\simeq} \underline{\hand{\Omega}}^1_{X/k}(x)$ when x is a rat/k point of a

\mspage{\hand{141}}
\tnote{25-141-1}{The live uncertainty translated as take has the French readings prendre/pirendre. The canceled final s after peut and the e after un are French-only typing or agreement changes with no English effect. The final abbreviation rat/k means rational over k. The change from Y in the setup to S in Proposition 5 is retained.}
